\section{Proofs for Section~\ref{sec:extensions}}

This appendix provides the algebra for the extensions in Section~\ref{sec:extensions}.

\subsection{Preliminaries}

Recall the baseline objects
\[
P_A=\frac{N_A(N_A-1)}{2},
\qquad
Q_{AB}=N_AN_B,
\qquad
\Delta m_A:=m_A^1-m_A^0>0,
\qquad
\Lambda_A:=(1+\beta)-\beta\delta(m_A^1+m_A^0).
\]

In the heterogeneous-participant extension, the expected surplus per realized local interaction is
\begin{equation}
\bar v_L(\sigma)
=
v_L^L(1-\sigma)^2
+
v_L^H\bigl[1-(1-\sigma)^2\bigr]
=
v_L^L+(2\sigma-\sigma^2)(v_L^H-v_L^L),
\label{eq:app7_vbarL_sigma}
\end{equation}
with
\begin{equation}
\bar v_L'(\sigma)=2(1-\sigma)(v_L^H-v_L^L)>0.
\label{eq:app7_vbarL_prime}
\end{equation}
Also, since $\bar v_L(\sigma)$ is a convex combination of $v_L^L$ and $v_L^H$,
\begin{equation}
\bar v_L(\sigma)>0
\qquad
\text{for all }\sigma\in[0,1].
\label{eq:app7_vbarL_positive}
\end{equation}

The expected surplus per effective cross-regional link is
\begin{equation}
\bar v_X(\sigma)
=
v_X^L(1-\sigma)+v_X^H\sigma
=
v_X^L+\sigma(v_X^H-v_X^L),
\label{eq:app7_vbarX_sigma}
\end{equation}
with
\begin{equation}
\bar v_X'(\sigma)=v_X^H-v_X^L>0
\label{eq:app7_vbarX_prime}
\end{equation}
and
\begin{equation}
\bar v_X(\sigma)>0
\qquad
\text{for all }\sigma\in[0,1].
\label{eq:app7_vbarX_positive}
\end{equation}

In the online-substitution extension, the no-hub matching rate is
\begin{equation}
m_A^0(o)=m_A^0+\phi o,
\qquad
\phi>0,
\label{eq:app7_m0o}
\end{equation}
with domain
\begin{equation}
0\le o<\bar o:=\frac{m_A^1-m_A^0}{\phi},
\label{eq:app7_o_domain}
\end{equation}
so that $0\le m_A^0(o)<m_A^1\le 1$.

\subsection{Proof of Proposition \ref{prop:heterogeneous_participants}}

Using \eqref{eq:vbarL_sigma}, the hub-alone welfare gain is
\[
\Delta_H(\sigma)
=
\bar v_L(\sigma)P_A\Delta m_A\Lambda_A-F.
\]
Differentiating with respect to $\sigma$,
\begin{align}
\frac{\partial \Delta_H(\sigma)}{\partial \sigma}
&=
\bar v_L'(\sigma)P_A\Delta m_A\Lambda_A.
\label{eq:app7_dDeltaH_dsigma_raw}
\end{align}
If $\Lambda_A>0$, then \eqref{eq:app7_vbarL_prime} implies
\[
\frac{\partial \Delta_H(\sigma)}{\partial \sigma}>0
\qquad
\text{for }\sigma\in[0,1),
\]
which proves \eqref{eq:dDeltaH_dsigma}.

Next, the optimal pipeline policy is
\[
b^*(\sigma)
=
\frac{\beta\eta \bar v_X(\sigma)Q_{AB}}{\kappa}.
\]
Differentiating,
\begin{align}
\frac{\partial b^*(\sigma)}{\partial \sigma}
&=
\frac{\beta\eta Q_{AB}}{\kappa}\bar v_X'(\sigma).
\label{eq:app7_dbstar_dsigma_raw}
\end{align}
Since all multiplicative terms are positive and $\bar v_X'(\sigma)>0$, we obtain
\[
\frac{\partial b^*(\sigma)}{\partial \sigma}>0,
\]
which proves \eqref{eq:dbstar_dsigma}.

For the renewal surplus,
\[
R_{AB}(N_A,N_B;\sigma)
=
\frac{\bigl(\beta\eta \bar v_X(\sigma)Q_{AB}\bigr)^2}{2\kappa}
=
\frac{\beta^2\eta^2Q_{AB}^2}{2\kappa}\,\bar v_X(\sigma)^2.
\]
Differentiating with respect to $\sigma$,
\begin{align}
\frac{\partial R_{AB}(N_A,N_B;\sigma)}{\partial \sigma}
&=
\frac{\beta^2\eta^2Q_{AB}^2}{2\kappa}\cdot 2\bar v_X(\sigma)\bar v_X'(\sigma)
\nonumber\\
&=
\frac{\beta^2\eta^2Q_{AB}^2}{\kappa}\,\bar v_X(\sigma)\bar v_X'(\sigma).
\label{eq:app7_dR_dsigma_raw}
\end{align}
By \eqref{eq:app7_vbarX_positive} and \eqref{eq:app7_vbarX_prime}, this derivative is strictly positive, proving \eqref{eq:dR_dsigma}.

Finally, consider the threshold statements. The short-run private-support condition in Section~\ref{sec:short_run} becomes
\[
F<\bar v_L(\sigma)P_A\Delta m_A.
\]
Thus, treating $N_A$ as continuous,
\begin{equation}
N_{SR}^*(\sigma)
=
\frac{1}{2}
\left(
1+\sqrt{1+\frac{8F}{\bar v_L(\sigma)\Delta m_A}}
\right).
\label{eq:app7_NSR_sigma}
\end{equation}
Differentiating,
\begin{align}
\frac{dN_{SR}^*(\sigma)}{d\sigma}
&=
-\frac{2F\,\bar v_L'(\sigma)}
{\Delta m_A\,\bar v_L(\sigma)^2
\sqrt{1+\frac{8F}{\bar v_L(\sigma)\Delta m_A}}}
<0.
\label{eq:app7_dNSR_dsigma}
\end{align}

Likewise, when $\Lambda_A>0$, the hub-efficiency threshold in Section~\ref{sec:hub_efficiency} becomes
\begin{equation}
N_H^*(\sigma)
=
\frac{1}{2}
\left(
1+\sqrt{1+\frac{8F}{\bar v_L(\sigma)\Delta m_A\Lambda_A}}
\right).
\label{eq:app7_NH_sigma}
\end{equation}
Differentiating,
\begin{align}
\frac{dN_H^*(\sigma)}{d\sigma}
&=
-\frac{2F\,\bar v_L'(\sigma)}
{\Delta m_A\Lambda_A\,\bar v_L(\sigma)^2
\sqrt{1+\frac{8F}{\bar v_L(\sigma)\Delta m_A\Lambda_A}}}
<0.
\label{eq:app7_dNH_dsigma}
\end{align}
Therefore both thresholds decrease with $\sigma$. This completes the proof of Proposition \ref{prop:heterogeneous_participants}.
\qed

\subsection{Proof of Proposition \ref{prop:online_substitution}}

In the online-substitution extension,
\[
m_A^0(o)=m_A^0+\phi o,
\qquad
0\le o<\bar o=\frac{m_A^1-m_A^0}{\phi}.
\]

The short-run private gain is
\[
\Pi_H(o)
=
v_LP_A\bigl(m_A^1-m_A^0(o)\bigr)-F
=
v_LP_A\bigl(m_A^1-m_A^0-\phi o\bigr)-F.
\]
Differentiating with respect to $o$,
\begin{align}
\frac{d\Pi_H(o)}{do}
&=
-\phi v_LP_A<0,
\label{eq:app7_dPiH_do_raw}
\end{align}
which proves \eqref{eq:dPiH_do}.

Next, the hub-alone welfare gain is
\[
\Delta_H(o)
=
v_LP_A\bigl[m_A^1-m_A^0(o)\bigr]
\Bigl[(1+\beta)-\beta\delta\bigl(m_A^1+m_A^0(o)\bigr)\Bigr]-F.
\]
For brevity, write $x:=m_A^0(o)$. Then
\[
\Delta_H(o)=v_LP_A(m_A^1-x)\Bigl[(1+\beta)-\beta\delta(m_A^1+x)\Bigr]-F.
\]
Differentiate first with respect to $x$:
\begin{align}
\frac{d}{dx}\Bigl[(m_A^1-x)\bigl((1+\beta)-\beta\delta(m_A^1+x)\bigr)\Bigr]
&=
-\Bigl[(1+\beta)-\beta\delta(m_A^1+x)\Bigr]
-(m_A^1-x)\beta\delta
\nonumber\\
&=
-(1+\beta)+\beta\delta(m_A^1+x)-\beta\delta(m_A^1-x)
\nonumber\\
&=
-(1+\beta)+2\beta\delta x.
\label{eq:app7_chain_x}
\end{align}
Since $dx/do=\phi$, the chain rule gives
\begin{align}
\frac{d\Delta_H(o)}{do}
&=
\phi v_LP_A\Bigl[-(1+\beta)+2\beta\delta m_A^0(o)\Bigr]
\nonumber\\
&=
-\phi v_LP_A\Bigl[(1+\beta)-2\beta\delta m_A^0(o)\Bigr].
\label{eq:app7_dDeltaH_do_raw}
\end{align}
This is exactly \eqref{eq:dDeltaH_do}.

To verify the sign, note that the domain restriction implies $m_A^0(o)<m_A^1$, and the regularity condition \eqref{eq:regularity} implies $\delta m_A^1<1$. Hence $\delta m_A^0(o)<1$ as well. Therefore
\[
2\beta\delta m_A^0(o)<2\beta.
\]
Since $\beta\in(0,1]$, we have $1+\beta\ge 2\beta$, with equality only when $\beta=1$. Thus
\[
(1+\beta)-2\beta\delta m_A^0(o)>(1+\beta)-2\beta=1-\beta\ge 0,
\]
where the strict inequality follows from $\delta m_A^0(o)<1$. Hence
\[
\frac{d\Delta_H(o)}{do}<0.
\]

Finally, conditional on a hub, welfare with pipeline policy is
\[
W(1,b)=W(1,0)+\beta\eta v_XQ_{AB}b-\frac{\kappa}{2}b^2.
\]
The planner therefore solves
\[
\max_{b\ge 0}
\left\{
W(1,0)+\beta\eta v_XQ_{AB}b-\frac{\kappa}{2}b^2
\right\}.
\]
The first-order condition is
\[
\beta\eta v_XQ_{AB}-\kappa b=0,
\]
so
\[
b^*(o)=\frac{\beta\eta v_XQ_{AB}}{\kappa}=b^*.
\]
Thus the optimal pipeline policy is unchanged by online substitution in this reduced form, which proves \eqref{eq:bstar_online}. This completes the proof of Proposition \ref{prop:online_substitution}.
\qed
